\section{Reduced schemes}
\label{section-reduced}

\begin{definition}
\label{definition-reduced}
Let $X$ be a scheme. We say $X$ is {\it reduced} if every local ring
$\mathcal{O}_{X, x}$ is reduced.
\end{definition}

\begin{lemma}
\label{lemma-reduced}
A scheme $X$ is reduced if and only if $\mathcal{O}_X(U)$
is a reduced ring for all $U \subset X$ open.
\end{lemma}

\begin{proof}
Assume that $X$ is reduced.
Let $f \in \mathcal{O}_X(U)$ be a section such that $f^n = 0$.
Then the image of $f$ in $\mathcal{O}_{U, u}$ is zero for
all $u \in U$. Hence $f$ is zero, see
Sheaves, Lemma \ref{sheaves-lemma-sheaf-subset-stalks}.
Conversely, assume that $\mathcal{O}_X(U)$ is reduced
for all opens $U$. Pick any nonzero element $f \in \mathcal{O}_{X, x}$.
Any representative $(U, f \in \mathcal{O}(U))$  of $f$ is nonzero and
hence not nilpotent. Hence $f$ is not nilpotent in $\mathcal{O}_{X, x}$.
\end{proof}

\begin{lemma}
\label{lemma-affine-reduced}
An affine scheme $\Spec(R)$ is reduced
if and only if $R$ is reduced.
\end{lemma}

\begin{proof}
The direct implication follows immediately from
Lemma \ref{lemma-reduced} above.
In the other direction it follows since any localization of
a reduced ring is reduced, and in particular the local rings
of a reduced ring are reduced.
\end{proof}

\begin{lemma}
\label{lemma-reduced-closed-subscheme}
Let $X$ be a scheme. Let $T \subset X$ be a closed subset.
There exists a unique closed subscheme $Z \subset X$ with
the following properties: (a) the underlying topological
space of $Z$ is equal to $T$, and (b) $Z$ is reduced.
\end{lemma}

\begin{proof}
Let $\mathcal{I} \subset \mathcal{O}_X$ be the sub presheaf
defined by the rule
$$
\mathcal{I}(U) = \{f \in \mathcal{O}_X(U) \mid
f(t) = 0\text{ for all }t \in T\cap U\}
$$
Here we use $f(t)$ to indicate the image of
$f$ in the residue field $\kappa(t)$ of $X$ at $t$.
Because of the local nature of the condition it is
clear that $\mathcal{I}$ is a sheaf of ideals. Moreover,
let $U = \Spec(R)$ be an affine open.
We may write $T \cap U = V(I)$ for a unique radical
ideal $I \subset R$. Given a prime $\mathfrak p \in V(I)$
corresponding to $t \in T \cap U$ and an element $f \in R$ we have
$f(t) = 0 \Leftrightarrow f \in \mathfrak p$.
Hence $\mathcal{I}(U) = \bigcap_{\mathfrak p \in V(I)} \mathfrak p
= I$ by Algebra, Lemma \ref{algebra-lemma-Zariski-topology}.
Moreover, for any standard open $D(g) \subset \Spec(R) = U$
we have $\mathcal{I}(D(g)) = I_g$ by the same reasoning.
Thus $\widetilde I$ and $\mathcal{I}|_U$ agree (as ideals)
on a basis of opens and hence are equal. Therefore
$\mathcal{I}$ is a quasi-coherent sheaf of ideals.

\medskip\noindent
At this point we may define $Z$ as the closed subspace
associated to the sheaf of ideals $\mathcal{I}$. For every
affine open $U = \Spec(R)$ of $X$ we see that
$Z \cap U = \Spec(R/I)$ where $I$ is a radical ideal and
hence $Z$ is reduced (by Lemma \ref{lemma-affine-reduced} above).
By construction the underlying closed subset of $Z$ is $T$.
Hence we have found a closed subscheme with properties (a) and (b).

\medskip\noindent
Let $Z' \subset X$ be a second closed subscheme with properties (a) and (b).
For every affine open $U = \Spec(R)$ of $X$ we see that
$Z' \cap U = \Spec(R/I')$ for some ideal $I' \subset R$.
By Lemma \ref{lemma-affine-reduced} the ring $R/I'$ is
reduced and hence $I'$ is radical. Since $V(I') = T \cap U = V(I)$
we deduced that $I = I'$ by
Algebra, Lemma \ref{algebra-lemma-Zariski-topology}.
Hence $Z'$ and $Z$ are defined by the same sheaf of ideals
and hence are equal.
\end{proof}

\begin{definition}
\label{definition-reduced-induced-scheme}
Let $X$ be a scheme. Let $Z \subset X$ be a closed subset.
A {\it scheme structure on $Z$} is given by a closed subscheme $Z'$ of
$X$ whose underlying set is equal to $Z$. We often say
``let $(Z, \mathcal{O}_Z)$ be a scheme structure on $Z$'' to
indicate this. The {\it reduced induced scheme structure}
on $Z$ is the one constructed in Lemma \ref{lemma-reduced-closed-subscheme}.
The {\it reduction $X_{red}$ of $X$} is the reduced induced scheme
structure on $X$ itself.
\end{definition}

\noindent
Often when we say ``let $Z \subset X$ be an irreducible component of $X$''
we think of $Z$ as a reduced closed subscheme of $X$ using the reduced induced
scheme structure.

\begin{remark}
\label{remark-reduced-induced-locally-closed}
Let $X$ be a scheme. Let $T \subset X$ be a locally closed subset.
In this situation we sometimes also use the phrase
``reduced induced scheme structure on $T$''. It refers
to the reduced induced scheme structure from
Definition \ref{definition-reduced-induced-scheme}
when we view $T$ as a closed subset of the open subscheme 
$X \setminus \partial T$ of $X$. Here
$\partial T = \overline{T} \setminus T$ is the ``boundary'' of $T$
in the topological space of $X$.
\end{remark}

\begin{lemma}
\label{lemma-map-into-reduction}
Let $X$ be a scheme.
Let $Z \subset X$ be a closed subscheme.
Let $Y$ be a reduced scheme.
A morphism $f : Y \to X$ factors through $Z$ if and only if
$f(Y) \subset Z$ (set theoretically). In particular, any
morphism $Y \to X$ factors as $Y \to X_{red} \to X$.
\end{lemma}

\begin{proof}
Assume $f(Y) \subset Z$ (set theoretically).
Let $\mathcal{I} \subset \mathcal{O}_X$ be the ideal sheaf of $Z$.
For any affine opens $U \subset X$, $\Spec(B) = V \subset Y$
with $f(V) \subset U$ and any $g \in \mathcal{I}(U)$
the pullback $b = f^\sharp(g) \in \Gamma(V, \mathcal{O}_Y) = B$
maps to zero in the residue field of any $y \in V$.
In other words $b \in \bigcap_{\mathfrak p \subset B} \mathfrak p$.
This implies $b = 0$ as $B$ is reduced (Lemma \ref{lemma-reduced}, and
Algebra, Lemma \ref{algebra-lemma-Zariski-topology}).
Hence $f$ factors through
$Z$ by Lemma \ref{lemma-characterize-closed-subspace}.
\end{proof}










